\chapter{The Sphere Theorem}
\label{chap:dc-sphere-theorem}

\begin{theorem}[Sphere Theorem]
  \label{thm:dc-ch13-1-1}
  \dcref{ch13:1.1}
  Let $M^n$ be a compact simply connected Riemannian manifold, whose sectional
  curvature $K$ satisfies

  $$
  0<hK_{\max}<K\leq K_{\max}.\qquad(1)
  $$

  Then if $h=1/4$, $M$ is homeomorphic to a sphere.

  The number $h$ is the pinching constant. Rescaling the metric gives the equivalent
  normalization

  $$
  0<h<K\leq 1.\qquad(1')
  $$

  In even dimension the non-strict bound

  $$
  0<1/4\leq K\leq1,\tag{2}
  $$
  admits simply connected examples not homeomorphic to a sphere in even dimensions.
  Under (2), the alternatives are $\operatorname{diam}(M)>\pi$ with $M$ homeomorphic
  to a sphere, or $\operatorname{diam}(M)=\pi$ with $M$ isometric to a symmetric
  space. In odd dimensions (2) still implies the conclusion. For $n=2,3$, positive
  sectional curvature suffices for a compact simply connected manifold to be
  homeomorphic to $S^n$.

  Unless otherwise stated, manifolds are complete and geodesics are normalized.
\end{theorem}
\begin{proof}
  \sketch
  Let $p,q\in M$ with $\operatorname{diam}M=d(p,q)$. Let $D_1$,
  $D_2$ be the subsets of $M$ formed by all geodesic segments $\overline{pm}$ and
  $\overline{qn}$ respectively, with $m,n$ from \cref{lem:dc-ch13-4-3}. By continuity of $m$, $D_1,D_2$ are closed; we claim $D_1\cup D_2=M$ and
  $\partial D_1=\partial D_2=D_1\cap D_2$. Indeed, for $r\in M$, by \cref{lem:dc-ch13-4-2} either
  $d(p,r)<\rho$ or $d(q,r)<\rho$; in the first case, since $d(p,C_m(p))\geq\pi>\rho$,
  there is a unique minimizing $\gamma$ through $p,r$ and a unique $m$ on $\gamma$ with
  $d(p,m)=d(q,m)<\rho$, and $r\in\overline{pm}=D_1$ if $d(p,r)<d(q,r)$, $r\in\partial D_1$
  if $d(p,r)=d(q,r)$ (then $r=m$), and $r\in D_2$ if $d(p,r)>d(q,r)$; so
  $r\in D_1\cup D_2$, and $r\in D_1\cap D_2$ forces $d(p,r)=d(q,r)$, $r=m=n$.

  Define $\varphi:S^n\to M^n$ as follows. Associate to a fixed $N\in S^n$ the point
  $p$, and to its antipode $S$ the point $q$. Choose a linear isometry
  $i:T_NS^n\to T_pM$. For each $e$ in the equator $E$ of $S^n$ relative to $N$, take
  the geodesic $\gamma(s)$ of $S^n$, $0\leq s\leq\pi$, with $\gamma(0)=N$,
  $\gamma(\pi/2)=e$; let $m$ be the point of \cref{lem:dc-ch13-4-3} on the geodesic of $M$ through
  $p$ with velocity $i\gamma'(0)$. Define

  $$
  \varphi(\gamma(s))=\exp_p\!\Big(s\tfrac{2}{\pi}d(p,m)\,(i\circ\gamma'(0))\Big),\quad 0<s\leq\tfrac{\pi}{2},
  $$

  $$
  \varphi(\gamma(s))=\exp_q\!\Big(\big(2-\tfrac{2s}{\pi}\big)d(q,m)\,V\Big),\quad\tfrac{\pi}{2}\leq s<\pi,
  $$

  where $V$ is the unit tangent at $q$ of the unique minimizing geodesic from $q$ to
  $m$. Then $\varphi$ maps the closed northern hemisphere bijectively onto $D_1$, the
  closed southern hemisphere bijectively onto $D_2$, and $E$ onto
  $\partial D_1=\partial D_2=D_1\cap D_2$. By uniqueness of $m,n$ (\cref{lem:dc-ch13-4-3}),
  $\varphi$ is continuous; since $M=D_1\cup D_2$, $\varphi$ is surjective; since
  $D_1\cap D_2=\partial D_1=\varphi(E)$, $\varphi$ is injective. As $S^n$ is compact,
  $\varphi$ is a homeomorphism.
\end{proof}

\section{The cut locus}

Let $M$ be complete, $p\in M$, and $\gamma:[0,\infty)\to M$ a normalized geodesic
with $\gamma(0)=p$. For small $t>0$, $d(\gamma(0),\gamma(t))=t$, and once a segment
ceases to be minimizing no later segment is minimizing. Thus the set of $t>0$ with
$d(\gamma(0),\gamma(t))=t$ is of the form $[0,t_0]$ or $[0,\infty)$. In the first
case $\gamma(t_0)$ is called the \emph{cut point of $p$ along $\gamma$}; in the second
no cut point exists. The \emph{cut locus of $p$}, denoted $C_m(p)$, is the union of the
cut points of $p$ along all geodesics starting from $p$. If $M$ is compact, a cut
point exists along every geodesic from $p$; the converse holds by
\cref{cor:dc-ch13-2-11}.

\begin{proposition}[Characterization of a cut point]
  \label{prop:dc-ch13-2-2}
  \uses{prop:dc-ch13-2-9}
  \dcref{ch13:2.2}
  Suppose that $\gamma(t_0)$ is the cut point of $p=\gamma(0)$ along $\gamma$. Then:
  (a) either $\gamma(t_0)$ is the first conjugate point of $\gamma(0)$ along
    $\gamma$,
  (b) or there exists a geodesic $\sigma\neq\gamma$ from $p$ to $\gamma(t_0)$ such
    that $\ell(\sigma)=\ell(\gamma)$.

  Conversely, if (a) or (b) holds, then there exists $\tilde t\in(0,t_0]$ such that
  $\gamma(\tilde t)$ is the cut point of $p$ along $\gamma$.
\end{proposition}
\begin{proof}
\sketch
  Take a sequence $\{t_0+\varepsilon_i\}$, $\varepsilon_i>0$,
  $\varepsilon_i\to 0$, and minimizing geodesics $\sigma_i$ joining $p$ to
  $\gamma(t_0+\varepsilon_i)$. The unit sphere in $T_pM$ being compact, pass to a
  subsequence with $\sigma_i'(0)\to\sigma'(0)$; by continuity $\sigma$ is a
  minimizing geodesic from $p$ to $\gamma(t_0)$, so $\ell(\sigma)=\ell(\gamma)$. If
  $\sigma\neq\gamma$, (b) holds. If $\sigma=\gamma$, we show $d\exp_p$ is singular at
  $t_0\gamma'(0)$: otherwise $\exp_p$ is a diffeomorphism on a neighborhood $U$ of
  $t_0\gamma'(0)$, and writing $\gamma(t_0+\varepsilon_j)=\sigma_j(t_0+\varepsilon_j')$
  with $\varepsilon_j'\leq\varepsilon_j$ gives
  $\exp_p(t_0+\varepsilon_j)\gamma'(0)=\exp_p(t_0+\varepsilon_j')\sigma_j'(0)$, hence
  $\gamma'(0)=\sigma_j'(0)$, contradicting the definition of $t_0$. Conversely, if
  (a) holds, since a geodesic does not minimize after the first conjugate point
  (\cref{cor:dc-ch11-2-9}), the cut point occurs at
  $\gamma(\tilde t)$, $\tilde t\leq t_0$. If (b) holds, joining $\sigma(t_0-\varepsilon)$
  to $\gamma(t_0+\varepsilon)$ by the unique minimizing geodesic $\tau$ inside a
  totally normal neighborhood of $\gamma(t_0)$ yields a curve of length strictly less
  than $t_0+\varepsilon$, so again the cut point occurs at $\gamma(\tilde t)$,
  $\tilde t\leq t_0$.
\end{proof}

\begin{example}[Cut locus of a sphere]
  \label{ex:dc-ch13-2-3}
  \dcref{ch13:2.3}
  If $M^n=S^n$ and $p\in S^n$, the cut locus of $p$ is its antipodal point. Here the
  cut locus of $p$ coincides with the locus of conjugate points to $p$.
\end{example}

\begin{example}[Cut and conjugate loci of real projective space]
  \label{ex:dc-ch13-2-4}
  \dcref{ch13:2.4}
  If $M^n=P^n(\mathbb{R})$, obtained by identifying antipodal points of $S^n$, the
  cut locus of $[p,-p]$ is the subset $P^{n-1}(\mathbb{R})\subset P^n(\mathbb{R})$
  obtained by identifying the antipodal points of the "equator" of $p$. The conjugate
  locus of $[p,-p]$ is $[p,-p]$ itself.
\end{example}

\begin{example}[Cut loci in flat torus and cylinder]
  \label{ex:dc-ch13-2-5}
  \dcref{ch13:2.5}
  On the flat torus $T$ obtained by identifying opposite sides of a quadrilateral
  $ABCD$, the cut locus of $A$ is the set formed by the medians of the segments $BA$
  and $BC$; the conjugate locus is empty. Similarly, the cut locus of a point $p$ of
  a cylinder in $\mathbb{R}^3$ is the "opposite generator" to that through $p$.
\end{example}

\begin{remark}[Cut and conjugate alternatives]
  \label{rem:dc-ch13-2-6}
  \uses{ex:dc-ch13-2-3, ex:dc-ch13-2-4, ex:dc-ch13-2-5, prop:dc-ch13-2-2}
  \dcref{ch13:2.6}
  In \cref{ex:dc-ch13-2-3} the situations (a) and (b) of \cref{prop:dc-ch13-2-2} occur simultaneously;
  in \cref{ex:dc-ch13-2-4} situation (b) occurs before (a); in \cref{ex:dc-ch13-2-5} only situation (b)
  occurs. For the ellipsoid, except for two directions, (b) occurs before (a).
\end{remark}

\begin{corollary}[Symmetry of cut points]
  \label{cor:dc-ch13-2-7}
  \uses{prop:dc-ch13-2-2}
  \dcref{ch13:2.7}
  If $q$ is a cut point of $p$ along $\gamma$, then $p$ is the cut point of $q$ along
  $-\gamma$; in particular $q\in C_m(p)$ if and only if $p\in C_m(q)$.
\end{corollary}
\begin{proof}
\sketch
  By \cref{prop:dc-ch13-2-2}, either $q$ is conjugate to $p$, or there is a geodesic
  $\sigma\neq\gamma$ from $p$ to $q$ with $\ell(\sigma)=\ell(\gamma)=d(p,q)$. In both
  cases the cut point of $q$ along $-\gamma$ does not occur before $p$; since
  $\ell(-\gamma)=d(p,q)$, $p$ is that cut point.
\end{proof}

\begin{corollary}[Injectivity radius]
  \label{cor:dc-ch13-2-8}
  \dcref{ch13:2.8}
  If $q\in M-C_m(p)$, there exists a unique minimizing geodesic joining $p$ to $q$.

  \cref{cor:dc-ch13-2-8} shows that $\exp_p$ is injective on the open ball $B_r(p)$ if and only
  if $r\leq d(p,C_m(p))$. For this reason

  $$
  i(M)=\inf_{p\in M}d(p,C_m(p))
  $$

  is called the \emph{injectivity radius of $M$}. Moreover, $M-C_m(p)$ is
  homeomorphic to an open Euclidean ball.
\end{corollary}

\begin{proposition}[Continuity of cut distance]
  \label{prop:dc-ch13-2-9}
  \dcref{ch13:2.9}
  Define $f:T_1M\to\mathbb{R}\cup\{\infty\}$ by $f(\gamma(0),\gamma'(0))=t_0$ if
  $\gamma(t_0)$ is the cut point of $\gamma(0)$ along $\gamma$, and
  $f(\gamma(0),\gamma'(0))=\infty$ if no cut point exists, where $\mathbb{R}\cup\{\infty\}$
  carries the topology with base the open intervals $(a,b)$ together with the sets
  $(a,\infty]=(a,\infty)\cup\{\infty\}$. Then $f$ is continuous.
\end{proposition}
\begin{proof}
\sketch
  Let $\gamma_i(0)\to\gamma(0)$, $\gamma_i'(0)\to\gamma'(0)$, with cut points
  $\gamma_i(t_o^i)$ and $\gamma(t_o)$. We show $\lim t_o^i=t_o$. First
  $\limsup t_o^i\leq t_o$: if $t_o<\infty$ and $\varepsilon>0$, only finitely many $j$
  can have $t_o+\varepsilon<t_o^j$, else $d(\gamma_j(0),\gamma_j(t_o+\varepsilon))=t_o+\varepsilon$
  gives in the limit $d(\gamma(0),\gamma(t_o+\varepsilon))=t_o+\varepsilon$,
  contradicting that $\gamma(t_o)$ is a cut point. Now let
  $\bar t=\liminf t_o^i\leq t_o$; suppose $\bar t<\infty$ and pass to a subsequence
  $t_o^j\to\bar t$. If for a subsequence the $\gamma_j(t_o^j)$ are conjugate to
  $\gamma_j(0)$, then (accumulation of conjugate points is conjugate) $\gamma(\bar t)$
  is conjugate to $\gamma(0)$, so $\bar t\geq t_o$. Otherwise, by \cref{prop:dc-ch13-2-2},
  geodesics $\sigma_j\neq\gamma_j$ exist with $\sigma_j(0)=\gamma_j(0)$,
  $\sigma_j(t_o^j)=\gamma_j(t_o^j)$, $\ell(\sigma_j)=\ell(\gamma_j)$; a subsequence
  $\sigma_j\to\sigma$ joins $\gamma(0)$ to $\gamma(\bar t)$. If $\sigma\neq\gamma$,
  \cref{prop:dc-ch13-2-2} gives $t_o\leq\bar t$; if $\sigma=\gamma$, a similar argument shows
  $\gamma(\bar t)$ conjugate to $\gamma(0)$, so $\bar t\geq t_o$.
\end{proof}

\begin{corollary}[Closedness of the cut locus]
  \label{cor:dc-ch13-2-10}
  \uses{prop:dc-ch13-2-9}
  \dcref{ch13:2.10}
  For all $p\in M$, $C_m(p)$ is closed; in particular, if $M$ is compact, $C_m(p)$ is
  compact.
\end{corollary}
\begin{proof}
\sketch
  One verifies $C_m(p)=\{\gamma(t);t=f(p,\gamma'(0))\}$, where $f$ is the
  function of \cref{prop:dc-ch13-2-9}. If $q$ is an accumulation point of $C_m(p)$, take
  $\gamma_j(t_j)\to q$ with $t_j=f(p,\gamma_j'(0))$ and $\gamma_j'(0)\to v$. With
  $\gamma(0)=p$, $\gamma'(0)=v$, continuity of $f$ gives
  $q=\exp_p(f(p,\gamma'(0))\gamma'(0))=\gamma(f(p,\gamma'(0)))\in C_m(p)$.
\end{proof}

\begin{corollary}[Cut point in every direction implies compactness]
  \label{cor:dc-ch13-2-11}
  \dcref{ch13:2.11}
  Suppose $M$ is complete and there exists $p\in M$ which has a cut point for every
  geodesic starting from $p$. Then $M$ is compact.
\end{corollary}
\begin{proof}
\sketch
  Since $M=\bigcup\{\gamma(t);t\leq f(p,\gamma'(0))\}$ over all geodesics
  $\gamma$ with $\gamma(0)=p$, and $f$ is continuous and (by hypothesis) bounded, $M$
  is bounded; the result follows from the Hopf–Rinow Theorem.
\end{proof}

\begin{proposition}[Geometry of a nearest cut point]
  \label{prop:dc-ch13-2-12}
  \uses{prop:dc-ch13-2-2}
  \dcref{ch13:2.12}
  Let $p\in M$ and suppose there exists $q\in C_m(p)$ realizing the distance from $p$
  to $C_m(p)$. Then:
  (a) either there exists a minimizing geodesic $\gamma$ from $p$ to $q$ along which
    $q$ is conjugate to $p$,
  (b) or there exist exactly two minimizing geodesics $\gamma$ and $\sigma$ from $p$
    to $q$, with $\gamma'(\ell)=-\sigma'(\ell)$, $\ell=d(p,q)$.
\end{proposition}
\begin{proof}
\sketch
  Let $\gamma$ be minimizing from $p$ to $q$. By \cref{prop:dc-ch13-2-2} either $q$
  is conjugate to $p$ along $\gamma$ — case (a) — or there is another minimizing
  $\sigma\neq\gamma$ with $\ell(\sigma)=\ell(\gamma)$. Suppose $q$ is not conjugate to
  $p$ along $\gamma$ or $\sigma$, and $\gamma'(\ell)\neq-\sigma'(\ell)$; we derive a
  contradiction (which also shows there can be only two such geodesics). Choose
  $V\in T_qM$ with $\langle V,\gamma'(\ell)\rangle<0$ and
  $\langle V,\sigma'(\ell)\rangle<0$. Since $q$ is not conjugate along $\gamma$, build
  a variation $\gamma_s$ of $\gamma$ with
  $\frac{d}{ds}\ell(\gamma_s)|_{s=0}=\langle V,\gamma'(\ell)\rangle<0$ by the formula
  for the first variation, and similarly $\sigma_s$ of $\sigma$ with
  $\frac{d}{ds}\ell(\sigma_s)|_{s=0}=\langle V,\sigma'(\ell)\rangle<0$ and
  $\sigma_s(\ell)=\tau(s)$. For small $s>0$, $\ell(\gamma_s)<\ell(\gamma)$ and
  $\ell(\sigma_s)<\ell(\sigma)$. If $\ell(\gamma_s)=\ell(\sigma_s)$, by \cref{prop:dc-ch13-2-2}
  $\tau(s)=\gamma_s(\ell)$ is a cut point with
  $d(p,\gamma_s(\ell))=\ell(\gamma_s)<d(p,C_m(p))$, a contradiction; if
  $\ell(\gamma_s)<\ell(\sigma_s)$, $\sigma_s$ is not minimizing so a cut point
  $\sigma_s(\tilde t)$, $\tilde t<\ell$, contradicts again; the case
  $\ell(\gamma_s)>\ell(\sigma_s)$ is similar.
\end{proof}

\begin{proposition}[Injectivity radius or a shortest closed geodesic]
  \label{prop:dc-ch13-2-13}
  \uses{prop:dc-ch13-2-9, prop:dc-ch10-2-4, prop:dc-ch13-2-12}
  \dcref{ch13:2.13}
  If the sectional curvature $K$ of a complete Riemannian manifold $M$ satisfies
  $0<K_{\min}\leq K\leq K_{\max}$, then:
  (a) $i(M)\geq\pi/\sqrt{K_{\max}}$, or
  (b) there exists a closed geodesic $\gamma$ in $M$, whose length is less than that
    of any other closed geodesic in $M$, with $i(M)=\tfrac{1}{2}\ell(\gamma)$.
\end{proposition}
\begin{proof}
\sketch
  By Bonnet--Myers, $M$ is compact. Since $T_1M$ is
  compact, \cref{prop:dc-ch13-2-9} gives $p\in M$ with $d(p,C_m(p))=\inf_{r\in M}d(r,C_m(r))$,
  and $C_m(p)$ compact gives $q\in C_m(p)$ realizing this distance. If $q$ is conjugate
  to $p$, then $d(p,q)\geq\pi/\sqrt{K_{\max}}$ by \cref{prop:dc-ch10-2-4} — case
  (a). If $q$ is not conjugate to $p$, by \cref{prop:dc-ch13-2-12} there are two minimizing
  geodesics $\mu,\sigma$ from $p$ to $q$ with $\mu'(\ell)=-\sigma'(\ell)$,
  $\ell=d(p,q)$. Since $q\in C_m(p)$ we have $p\in C_m(q)$ and $p$ realizes the
  distance from $q$ to $C_m(q)$, so $\mu'(0)=-\sigma'(0)$; thus $\mu$ and $\sigma$ form
  a closed geodesic $\gamma$ satisfying (b).
\end{proof}

\section{The estimate of the injectivity radius}

\begin{proposition}[Lower bound for the injectivity radius]
  \label{prop:dc-ch13-3-1}
  \dcref{ch13:3.1}
  Let $M^n$, $n\geq 3$, be a simply connected, compact Riemannian manifold with
  $1/4<K\leq 1$. Then $i(M)\geq\pi$.
\end{proposition}
\begin{proof}
  \sketch
  Suppose $i(M)<\pi$. By \cref{prop:dc-ch13-2-13} there is a
  closed geodesic $\gamma$ of length $\ell=\ell(\gamma)<2\pi$. By \cref{cor:dc-ch11-2-10} the conjugate points to $\gamma(0)$ along $\gamma$ are discrete; choose
  $\varepsilon>0$ with: (1) $\gamma(\ell-\varepsilon)$ not conjugate to $p=\gamma(0)$;
  (2) $\exp_p$ a diffeomorphism on $B_{2\varepsilon}(p)$;
  (3) $3\varepsilon<2\pi-\pi/\sqrt{\bar K}$, $\bar K=\inf_M K$;
  (4) $3\varepsilon<2\pi-\ell$; (5) $5\varepsilon<2\pi$. By Sard's Theorem pick a
  regular value $q\in B_\varepsilon(\gamma(\ell-\varepsilon))$ of $\exp_p$, with
  $q=\gamma_1(t)$ a geodesic from $p$, $3\varepsilon<\ell(\gamma_1)<\ell$. A
  minimizing geodesic $\gamma_o$ from $p$ to $q$ has
  $\ell(\gamma_o)\leq d(p,\gamma(\ell-\varepsilon))+d(\gamma(\ell-\varepsilon),q)\leq 2\varepsilon$,
  so $\gamma_o\neq\gamma_1$.

  Consider $\Omega_{p,q}^c$ (\cref{rem:dc-ch11-2-12}) and its finite-dimensional
  approximation $B$. Since $q$ is a regular value, all critical points in
  $\Omega_{p,q}^c$ are non-degenerate. As $M$ is simply connected, there is a homotopy
  $\gamma_s$ between $\gamma_o$ and $\gamma_1$ fixing $p,q$; deforming into $B$ and
  applying the index-zero/one variant of \cref{lem:dc-ch13-3-3}, for all
  $\delta>0$ the larger-length curve
  $\bar\gamma$ of the deformed $\bar\gamma_s$ satisfies $\ell(\bar\gamma)<a+\delta$,
  $a=\max\{\ell(\gamma_o),\ell(\gamma_1),\ell(\sigma)\}$, with $\sigma$ the largest
  geodesic of index $<2$. Put $\delta=\varepsilon$. Here $\ell(\gamma_o)\leq2\varepsilon$
  and $\ell(\gamma_1)<\ell<2\pi-3\varepsilon$ by (4). To estimate $\ell(\sigma)$, apply
  \cref{lem:dc-ch13-3-2} with $\bar M$ a sphere of curvature $\bar K=\inf_M K>1/4$, getting
  $\operatorname{index}\bar\sigma\leq\operatorname{index}\sigma<2$. Since $\bar\sigma$
  is a geodesic on $S^n$, the Morse Index Theorem gives: if
  $\ell(\sigma)>\pi/\sqrt{\bar K}$ then $\operatorname{index}\bar\sigma\geq2$;
  therefore by (3) $\ell(\sigma)\leq\pi/\sqrt{\bar K}<2\pi-3\varepsilon$. Hence by (5)
  $\ell(\bar\gamma)<a+\varepsilon<2\pi-3\varepsilon+\varepsilon=2\pi-2\varepsilon$.
  On the other hand, the closed-geodesic estimate gives in the
  homotopy a curve $\bar\gamma_{s_o}$ with $\ell(\gamma_o)+\ell(\bar\gamma_{s_o})\geq2\pi$,
  so $\ell(\bar\gamma)\geq2\pi-\ell(\gamma_o)\geq2\pi-2\varepsilon$, contradicting
  $\ell(\bar\gamma)<2\pi-2\varepsilon$. Thus $i(M)<\pi$ is untenable.
\end{proof}

\begin{lemma}[Index comparison]
  \label{lem:dc-ch13-3-2}
  \dcref{ch13:3.2}
  Let $M$ and $\bar M$ be Riemannian manifolds of the same dimension with sectional
  curvatures satisfying $\sup\bar K\leq\inf K$. Consider a geodesic
  $\sigma:[0,\ell]\to M$, $\sigma(0)=p$, fix $\bar p\in\bar M$, let
  $i:T_pM\to T_{\bar p}\bar M$ be a linear isometry, and set
  $\bar\sigma(t)=\exp_{\bar p}t(i\sigma'(0))$. Then
  $\operatorname{index}\sigma\geq\operatorname{index}\bar\sigma$.
\end{lemma}
\begin{proof}
\sketch
  For a piecewise differentiable field $\bar W$ along $\bar\sigma$, define
  $W$ along $\sigma$ by $W(t)=P_t\circ i^{-1}\circ\bar P_t^{-1}(\bar W(t))$ (parallel
  transports). Then $\langle W,\sigma'\rangle=\langle\bar W,\bar\sigma'\rangle$,
  $|W(t)|=|\bar W(t)|$, $\langle W',W'\rangle=\langle\bar W',\bar W'\rangle$. Since
  $\sup\bar K\leq\inf K$, $I(W,W)\leq I(\bar W,\bar W)$; so $I(\bar W,\bar W)<0$ forces
  $I(W,W)<0$.
\end{proof}

For the next lemma we need the following facts. Let $f:M^n\to\mathbb{R}$ be a
differentiable function on a differentiable manifold $M$. A point $p\in M$ is a
critical point of $f$ if $df(p)=0$; $f(p)$ is then called a critical value of $f$.
If $p$ is a critical point and $(x_1,\dots,x_n)$ is a coordinate system on $M$
around $p$, the hessian of $f$, defined by the matrix
$(\partial^2f/\partial x_i\partial x_j)(p)$, represents a symmetric bilinear form
on $T_pM$ that does not depend on the chosen system of coordinates. We say that
the critical point is non-degenerate if the determinant of this matrix is
different from zero. Such a critical point is isolated, and it is possible to
choose a coordinate system $(x_1,\dots,x_{n-\lambda},y_1,\dots,y_\lambda)$ in such
a way that, in the coordinate neighborhood considered,

$$
f(x_1,\dots,x_{n-\lambda},y_1,\dots,y_\lambda)
=f(p)+x_1^2+\cdots+x_{n-\lambda}^2-y_1^2-\cdots-y_\lambda^2.
$$

The integer $\lambda$ is the index of the non-degenerate critical point $p$.

\begin{lemma}[Morse-theoretic deformation]
  \label{lem:dc-ch13-3-3}
  \uses{prop:dc-ch13-3-1, prop:dc-ch13-2-13, prop:dc-ch13-2-12, lem:dc-ch13-3-2, cor:dc-ch13-2-10}
  \dcref{ch13:3.3}
  Let $M^n$ be a differentiable manifold and $f:M^n\to\mathbb{R}$ a function whose
  critical points are all non-degenerate. Given $p,q\in M$ and a differentiable curve
  $\gamma:[0,1]\to M$ with $\gamma(0)=p$, $\gamma(1)=q$, let $a=\max\{f(p),f(q)\}$,
  $M^a=\{x\in M;f(x)\leq a\}$, and $b=\max_{t\in[0,1]}(f\circ\gamma(t))$. Assume
  $f^{-1}([a,b])$ is compact and contains no critical points of index zero or one.
  Then for all $\delta>0$, $\gamma$ is homotopic, keeping endpoints fixed, to a curve
  $\bar\gamma$ with $\bar\gamma([0,1])\subset M^{a+\delta}$.
\end{lemma}
\begin{proof}
\sketch
  If $f^{-1}([a_1,a_2])$ is compact and
  critical-point-free, the gradient flow of $f$ gives a deformation retract
  $M^{a_2}\to M^{a_1}$. Let $p_1,\dots,p_k$ be the critical points in $M^b-M^a$ with
  distinct values $b\geq c_1>\dots>c_k>a$. At a critical value $c_1$ with index
  $\lambda\geq 2$, write $f=c_1+x_1^2+\dots+x_{n-\lambda}^2-y_1^2-\dots-y_\lambda^2$ on
  a neighborhood $H$, retract $M^{c_1+\varepsilon}\to M^{c_1-\varepsilon}\cup H$, and
  let $L=\{y_1=\dots=y_\lambda=0\}$. Since
  $\dim\bar\gamma_1'+\dim L\leq 1+n-2<\dim M$, transversality lets one perturb
  $\bar\gamma_1$ off $L$ within $H$ fixing endpoints; the gradient flow then retracts
  into $M^{c_1-\varepsilon}$. Proceeding inductively yields
  $\bar\gamma\subset M^{a+\delta}$.

  If $f^{-1}([a,b])$ contains critical points of index zero or one and $c$ is the
  largest value among them, the same argument gives
  $\bar\gamma([0,1])\subset M^{c+\delta}$.
\end{proof}

\begin{proposition}[Even-dimensional injectivity-radius estimate]
  \label{prop:dc-ch13-3-4}
  \uses{prop:dc-ch13-2-12, ex:dc-ch13-2-4}
  \dcref{ch13:3.4}
  If the sectional curvature $K$ of a compact orientable even-dimensional Riemannian
  manifold $M^n$ satisfies $0<K\leq 1$, then $i(M)\geq\pi$.
\end{proposition}
\begin{proof}
\sketch
  By compactness there are $p,q\in M$ with $q\in C_m(p)$ and $d(p,q)=i(M)$.
  Suppose $d(p,q)<\pi$. If $p$ is conjugate to $q$, then $d(p,q)\geq\pi$ by
  \cref{prop:dc-ch10-2-4}, a contradiction. So $p$ is not conjugate to $q$; by
  \cref{prop:dc-ch13-2-12,prop:dc-ch13-2-13} there is a closed geodesic $\gamma$ through $p=\gamma(0)$
  and $q$ with $\ell(\gamma)<2\pi$. Parallel transport along $\gamma$ leaves invariant
  a vector $V$ orthogonal to $\gamma$; computing the second
  variation $E_V''(0)$ for this $V$, it is strictly negative (here even dimension is
  crucial). Hence there is a variation through regular closed curves $\gamma_s$ with
  $\ell(\gamma_s)<\ell(\gamma)$ for $s\neq0$. Let $q_s$ be a point of $\gamma_s$ at
  maximum distance from $\gamma_s(0)$; the minimizing geodesic $\sigma_s$ from
  $q_s=\sigma_s(0)$ to $\gamma_s(0)$ has $\sigma_s(0)\to q$, and an accumulation point
  $w$ of $\sigma_s'(0)$ gives a minimizing geodesic $\sigma(t)=\exp_q tw$ from $q$ to
  $p$. By the first-variation formula $\sigma_s'(0)\perp\gamma_s'$ at $q_s$, so
  $\sigma'(0)\perp\gamma'$ at $q$. This yields three minimizing geodesics from $p$ to
  $q$; since $p$ is not conjugate to $q$, this contradicts \cref{prop:dc-ch13-2-12}.
\end{proof}

\begin{remark}[Orientability in even dimensions]
  \label{rem:dc-ch13-3-5}
  \uses{prop:dc-ch13-3-4}
  \dcref{ch13:3.5}
  By Synge's theorem the manifolds satisfying the hypotheses of
  \cref{prop:dc-ch13-3-4} are simply connected. Conversely a simply connected manifold is
  orientable (otherwise its orientation double cover would contradict simple
  connectivity). Hence the orientability hypothesis in
  \cref{prop:dc-ch13-3-4} is equivalent to assuming $M$ simply connected.
\end{remark}

\section{The Sphere Theorem}

In what follows, whether $\dim M$ is even or odd plays no role.

\begin{lemma}[Initial direction to a diameter endpoint]
  \label{lem:dc-ch13-4-1}
  \dcref{ch13:4.1}
  Let $M$ be a compact Riemannian manifold and $p,q\in M$ with
  $d(p,q)=\operatorname{diam}M$. Then for all $W\in T_pM$ there exists a minimizing
  geodesic $\gamma$ from $p=\gamma(0)$ to $q$ with $\langle\gamma'(0),W\rangle\geq0$.
\end{lemma}
\begin{proof}
\sketch
  Let $\lambda(t)=\exp_p tW$ and let $\gamma_t$ be a minimizing geodesic with
  $\gamma_t(0)=\lambda(t)$, $\gamma_t(\ell(\gamma_t))=q$. If for each $n$ there is
  $t_n\in[0,1/n]$ with $\langle\gamma_{t_n}'(0),\lambda'(t_n)\rangle\geq0$, a
  subsequence $\gamma_{t_n}\to\gamma$ gives
  $0\leq\langle\gamma'(0),\lambda'(0)\rangle=\langle\gamma'(0),W\rangle$. Otherwise
  some $n$ has $\langle\gamma_t'(0),\lambda'(t)\rangle<0$ for all $t\in[0,1/n]$. Taking
  a totally normal neighborhood $U$ of $\lambda(t)$, $q_o\in U$ on $\gamma_t$, and
  $\sigma_s$ minimizing from $q_o$ to $\lambda(s)$, the first variation of energy gives
  $\tfrac{1}{2}\frac{d}{ds}E(\sigma_s)|_{s=t}=-\langle\gamma_t'(0),\lambda'(t)\rangle>0$,
  so for $s<t$, $d(q_o,\lambda(s))<d(q_o,\lambda(t))$, hence
  $d(q,\lambda(s))<d(q,\lambda(t))$; but $p$ is at maximum distance from $q$, so
  $d(q,\lambda(t))\leq d(q,\lambda(0))$ — a contradiction.
\end{proof}

\begin{lemma}[Covering by two balls]
  \label{lem:dc-ch13-4-2}
  \uses{ex:dc-ch13-2-5}
  \dcref{ch13:4.2}
  Let $M^n$ be a compact, simply connected Riemannian manifold whose sectional
  curvature satisfies $\tfrac14<\delta\leq K\leq 1$, and let $p,q\in M$ with
  $d(p,q)=\operatorname{diam}(M)$. Then

  $$
  M=B_\rho(p)\cup B_\rho(q),
  $$

  where $B_r(p)$ denotes the open geodesic ball of radius $r$ and center $p$, and
  $\rho$ is such that $\pi/(2\sqrt\delta)<\rho<\pi$.
\end{lemma}
\begin{proof}
\sketch
  By \cref{prop:dc-ch13-3-1,prop:dc-ch13-3-4}, $B_\rho(p)$ contains no
  points of $C_m(p)$, so it is diffeomorphic via $\exp$ to a Euclidean ball, and
  likewise $B_\rho(q)$. (Here $i(M)\geq\pi$ enters crucially.) Suppose some $r\in M$
  has $d(p,r)\geq\rho$ and $d(q,r)\geq\rho$, with $d(p,r)\geq d(q,r)\geq\rho$. A
  minimizing geodesic from $q$ to $r$ meets $\partial B_\rho(q)$ at a point
  $q'\notin B_\rho(p)$ (else $d(r,q')>d(r,B_\rho(p))\geq d(r,B_\rho(q))=d(r,q')$,
  absurd). By Bonnet–Myers, $\operatorname{diam}(M)\leq\pi/\sqrt\delta<2\rho$, so the
  intersection $q''$ of the minimizing geodesic from $q$ to $p$ with $\partial B_\rho(q)$
  satisfies $d(p,q'')=d(p,q)-d(q,q'')<2\rho-\rho=\rho$, i.e. $q''\in B_\rho(p)$. As
  $\partial B_\rho(q)$ is path connected, $\partial B_\rho(p)\cap\partial B_\rho(q)\neq\emptyset$,
  giving $r_o$ with $d(r_o,p)=d(r_o,q)=\rho$. Take $\lambda$ minimizing from $p$ to
  $r_o$; by \cref{lem:dc-ch13-4-1} there is $\gamma$ minimizing from $p$ to $q$ with
  $\langle\gamma'(0),\lambda'(0)\rangle\geq0$. Let $s\in\gamma$ with $d(p,s)=\rho$.
  Applying Rauch comparison (\cref{prop:dc-ch10-2-5}), comparing $M^n$ with a
  sphere $S^n$ of curvature $\delta$: since the angle
  $\sphericalangle r_o p s\leq\pi/2$, $d(r_o,s)\leq\pi/(2\sqrt\delta)$. As
  $d(r_o,p)=d(r_o,q)=\rho$ and some $s$ has $d(r_o,s)<\rho$, the distance from $r_o$ to
  $\gamma$ is realized by an interior point $s_o$, with the minimizing geodesic from
  $r_o$ to $s_o$ orthogonal to $\gamma$ and
  $d(r_o,\gamma)=d(r_o,s_o)\leq\pi/(2\sqrt\delta)$. Since $d(p,q)\leq\pi/\sqrt\delta$,
  either $d(p,s_o)\leq\pi/(2\sqrt\delta)$ or $d(q,s_o)\leq\pi/(2\sqrt\delta)$; in the
  former case, since $d(r_o,s_o)\leq\pi/(2\sqrt\delta)$ and
  $\sphericalangle p s_o r_o=\pi/2$, Rauch gives $d(p,r_o)\leq\pi/(2\sqrt\delta)<\rho$,
  contradicting $d(p,r_o)=\rho$ (the other case is analogous).
\end{proof}

\begin{lemma}[Equator between two diameter endpoints]
  \label{lem:dc-ch13-4-3}
  \uses{lem:dc-ch13-4-2}
  \dcref{ch13:4.3}
  Under the conditions of \cref{lem:dc-ch13-4-2}, on each geodesic of length $\rho$ starting from
  $p$ there exists a unique point $m$ such that $d(p,m)=d(q,m)<\rho$. Similarly, on
  each geodesic of length $\rho$ starting from $q$ there exists a unique point $n$
  equidistant from $p$ and $q$.
\end{lemma}
\begin{proof}
\sketch
  For a geodesic $\gamma$ with $\gamma(0)=p$, set
  $f(s)=d(q,\gamma(s))-d(p,\gamma(s))$. Then $f$ is continuous and $f(0)=d(q,p)>0$. If
  $\gamma(s_o)$ is the cut point of $p$ along $\gamma$,
  \cref{prop:dc-ch13-3-1,prop:dc-ch13-3-4} gives
  $d(p,\gamma(s_o))\geq\pi>\rho$, and
  \cref{lem:dc-ch13-4-2} gives $d(q,\gamma(s_o))<\rho$, so $f(s_o)<0$. Hence $f(s_1)=0$ for some
  $s_1\in(0,s_o)$, giving $m=\gamma(s_1)$. For uniqueness, suppose $m_1\neq m_2$ both
  equidistant with $m_1$ between $p$ and $m_2$; then
  $d(q,m_2)=d(p,m_2)=d(p,m_1)+d(m_1,m_2)=d(q,m_1)+d(m_1,m_2)$, so the unique minimizing
  geodesic $\sigma_1$ from $q$ to $m_1$ coincides with $\gamma$, forcing $q\in\gamma$
  and then $p=q$, a contradiction.
\end{proof}

\begin{remark}[Continuity and sphere gluing]
  \label{rem:dc-ch13-4-4}
  \uses{thm:dc-ch13-1-1, lem:dc-ch13-4-2, lem:dc-ch13-4-3}
  \dcref{ch13:4.4}
  By its uniqueness, the point $m$ depends continuously on the (initial direction of
  the) geodesic which contains it.
\end{remark}

\begin{remark}[Remark 2.1]
  \label{rem:dc-ch13-2-1}
  \dcref{ch13:2.1}
  The cut locus, for surfaces, was introduced as the "ligne de partage" by
  H. Poincaré in 1905 [Po]. For Riemannian manifolds it was introduced by
  J. H. C. Whitehead in 1935 [Wh], who proved several of the properties presented
  here (pp. 700–704 of [Wh]). Interest revived in 1959 when Klingenberg [Kℓ 1]
  showed its usefulness in improving the "pinching" of the sphere theorem.
\end{remark}
